%!TEX root = ../../main.tex
\section{데이터 역할}
\label{sec:data-roles}

역할은 경기 단위로 배정된다. 15.14의 경기는 학습에 쓰고, 15.15의 경기는 보정과 선정에 쓰며, 15.16의 경기는 학습이 끝난 뒤 점수만 내는 평가에 쓴다. 한 경기에서 나온 교전은 모두 같은 역할에 들어가므로 같은 경기가 학습과 평가에 함께 나타나지 않는다. 이 배정은 두 모형에 서로 다른 단위로 적용된다. 승률 모형은 경기 안의 한 시각마다, 교전의 결과가 오를지를 예측하는 모형은 교전마다 나뉜다.

%% table 객체: tab:roles-evaluator — 출처 A_MLP_expanded_evaluator_meta_20260919.json meta.census; TEST_BAND_LEDGER
\begin{table}[!t]
  \centering
  \caption{승률 평가기 $\Vhat$의 데이터 역할 (시간상태 단위)}
  \label{tab:roles-evaluator}
  \small
  \begin{tabular}{@{}llrrl@{}}
    \toprule
    역할 & 패치 & 시간상태 행 & 경기 & 용도 \\
    \midrule
    TRAIN 적합 & 15.14 & \numVFit{} & --- & 가중치 \\
    TRAIN 정지 & 15.14 & \numVStop{} & --- & 조기 종료 (경기 단위 \holdoutV{}\% 보류, 시드 \bootSeed{}) \\
    V\_CAL & 15.15 & \numVCal{} & --- & 양의 기울기 시그모이드 $g$ \\
    V\_SELECT & 15.15 & \numVSel{} & --- & 시간 균형 손실 $L_{\mathrm{time}}$에 의한 선정 \\
    TEST & 15.16 & \numVTest{} & \numVTestMatches{} & 봉인 원장 \\
    \bottomrule
  \end{tabular}
  \tabnote{학습 합계 \numVTrain{} = 적합 \numVFit{} + 조기 종료용 \numVStop{}. ``---'' = census에 경기 수가 기록되지 않음.}
\end{table}

%% table 객체: tab:roles-engagement — 교전 단위 역할. 코드명·기호는 본문 말로 바꿈.
\begin{table}[!t]
  \centering
  \caption{한타와 소규모 교전을 학습과 평가에 나눈 수}
  \label{tab:roles-engagement}
  \footnotesize
  \begin{tabular}{@{}llccrr@{}}
    \toprule
    집단 & 쓰는 일 & 패치 & 지역 & 교전 & 경기 \\
    \midrule
    한타 & 학습 & 15.14 & 한국 & \numTTrain{} & \numTTrainMatches{} \\
    한타 & 보정 & 15.15 & 한국 & \numTQcal{} & \numTQcalMatches{} \\
    한타 & 선정 & 15.15 & 한국 & \numTQsel{} & \numTQselMatches{} \\
    한타 & 승률 모형의 보정 & 15.15 & 한국 & \numTVcal{} & \numTVcalMatches{} \\
    한타 & 승률 모형의 선정 & 15.15 & 한국 & \numTVsel{} & \numTVselMatches{} \\
    한타 & 평가 & 15.16 & 한국 & \numTTest{} & \numTTestMatches{} \\
    한타 & 평가의 일부 & 15.16 & 한국 & \numTBforty{} & \numTBfortyMatches{} \\
    한타 & 점수만 & 16.13 & 한국 & \numTExtKR{} & \numTExtKRMatches{} / \numExtKRCollected{} \\
    한타 & 점수만 & 16.13 & 북미 & \numTExtNA{} & \numTExtNAMatches{} / \numExtNACollected{} \\
    한타 & 보고만 & 16.15 & 한국 & \numTExtKRb{} & \numTExtKRbMatches{} / \numExtKRbCollected{} \\
    한타 & 보고만 & 16.14 & 한국 & \numTExtPilot{} & \numTExtPilotMatches{} / \numExtPilotCollected{} \\
    \midrule
    소규모 교전 & 학습 & 15.14 & 한국 & \numSTrain{} & \numSTrainMatches{} \\
    소규모 교전 & 보정 & 15.15 & 한국 & \numSQcal{} & \numSQcalMatches{} \\
    소규모 교전 & 선정 & 15.15 & 한국 & \numSQsel{} & \numSQselMatches{} \\
    소규모 교전 & 평가 & 15.16 & 한국 & \numSTest{} & \numSTestMatches{} \\
    소규모 교전 & 평가의 일부 & 15.16 & 한국 & \numSBforty{} & \numSBfortyMatches{} \\
    소규모 교전 & 점수만 & 16.13 & 한국 & \numSExtKR{} & --- / \numExtKRCollected{} \\
    소규모 교전 & 점수만 & 16.13 & 북미 & \numSExtNA{} & --- / \numExtNACollected{} \\
    소규모 교전 & 보고만 & 16.15 & 한국 & \numSExtKRb{} & --- / \numExtKRbCollected{} \\
    소규모 교전 & 보고만 & 16.14 & 한국 & \numSExtPilot{} & --- / \numExtPilotCollected{} \\
    \bottomrule
  \end{tabular}
\end{table}


표 \ref{tab:roles-evaluator}는 승률 모형 쪽이다. 여기서 세는 단위인 상태 하나는 경기 안의 한 시각에서 다시 만든 기록이고, 한 경기가 여러 상태를 낸다. 학습에 쓴 상태는 가중치를 맞춘 \numVFit{}과 학습을 멈출 때를 정한 \numVStop{}을 더한 \numVTrain{}이다. 경기 수가 ---인 행은 그 수가 기록에 남아 있지 않다는 뜻이다. 보정과 선정이 어떤 함수와 어떤 점수를 쓰는지는 제 \ref{sec:value-evaluator} 절에 있다.

표 \ref{tab:roles-engagement}는 교전 쪽이다. 집단 칸의 한타와 소규모 교전은 같은 교전을 공유하지 않는 별개의 집단이고, 지역 칸은 그 경기가 치러진 서버이다. 15.16까지는 모두 한국 서버이며, 북미 서버는 나중에 따로 모은 경기에만 나온다. 쓰는 일 칸에 ``승률 모형의''가 붙은 두 행은 승률 모형을 맞추고 고르는 데만 쓰고 교전 결과의 예측에는 넣지 않는다. 그 표시가 없는 보정 행과 선정 행이 교전 결과를 예측하는 모형의 몫이다. 15.15의 한타 \numTVal{}건은 이 네 행으로 모두 나뉜다.

한타에서는 보정 행에서 확률을 다시 맞추는 후보를 만들고, 선정 행에서 그 후보와 설정을 고른다. 소규모 교전은 보고하는 모형에서 확률을 다시 맞추지 않으므로 그 보정 행을 모형을 고르는 데 쓰지 않고, 선정 행에서 설정만 고른다. 설정이 무엇인지는 제 \ref{ch:prediction} 장에 있다.

나머지 세 종류의 행은 다음과 같이 읽는다. 평가의 일부는 평가 교전 가운데 시작 승률이 어느 쪽으로도 크게 기울지 않은 교전이며, 그 범위는 제 \ref{ch:value} 장에서 정한다. 이미 평가 행에 들어 있으므로 평가 행과 더하지 않는다. 점수만인 행은 모형을 다시 학습하지 않고 점수만 계산하는 경기이고, 보고만인 행은 주 결과에 넣지 않는다. 16.14는 제 \ref{sec:data-corpus} 절에서 파일럿으로 모았다고 한 경기이다. 경기 칸에 빗금이 있는 행은 앞이 그 집단에 교전을 남긴 경기, 뒤가 모은 경기 전체이다. 모은 경기는 한타와 소규모 교전이 같은 경기이므로 빗금 뒤의 수도 같다. 소규모 교전의 바깥 표본에서 빗금 앞이 ---인 것은, 그 가운데 소규모 교전을 남긴 경기가 몇인지가 기록에 남아 있지 않다는 뜻이다.

학습에 쓴 패치의 교전은 그 경기를 빼고 다시 만든 승률 모형이 결과를 매긴다. 제 \ref{sec:intro-approach} 절에서 말한 대로, 학습에 쓴 경기의 교전을 그 경기가 들어 있는 모형으로 매기지 않기 위해서이다. 한타와 소규모 교전은 같은 방식으로 다시 만든 모형을 쓴다. 보정, 선정, 평가, 그리고 나중에 모은 경기의 교전은 15.14로 학습한 뒤 고정한 승률 모형이 매긴다. 그 모형의 학습 경기는 15.14뿐이므로 그 뒤 패치의 경기와 겹치지 않는다.

두 집단을 이어 붙여 하나로 학습하는 방식은 제 \ref{ch:prediction} 장에서 따로 견주고, 점수를 내는 평가는 집단마다 둔다.
